% sections/methods.tex
% No internal phase labels or model filenames in the main text.
% Experiments described by what they test, not by planning labels.
% All run IDs and filenames go exclusively in Data Availability.

\subsection{Model Overview}

The GATA1/PU.1 commitment circuit was implemented as a stochastic model in the
Signal Hierarchical Petri Net formalism~\cite{shypn2026} (described in
Section~\ref{sec:shpn}). The model integrates nine biochemical subsystems: EPO
and GCSF receptor binding, internalisation, and recycling; GATA1 and PU.1 mRNA
transcription, splicing, and nuclear export; protein translation, nuclear import,
and phosphorylation; ubiquitin-mediated PU.1 degradation; and the ATP/GTP energy
nucleotide cycles that power these transitions. Each subsystem is represented by
one or more stochastic transitions whose firing rates combine mass-action kinetics
with Michaelis--Menten saturation and a shared pH-dependent inhibition term
governing cross-repression between GATA1 and PU.1 (equation~\ref{eq:pHhill}).

\subsection{Single-Cell and Population-Scale Regimes}

All stochastic simulations were run with initial molecular concentrations scaled
to one-hundredth of their standard values unless otherwise stated. At this scaling
the system operates in the physiological single-cell copy-number regime, where
stochastic fluctuations dominate and individual replicates yield heterogeneous
fate outcomes. Population-scale behaviour was verified by running a subset of
conditions at standard (unscaled) initial conditions, confirming that variance
collapsed and all replicates converged to the erythroid attractor.

\subsection{Experiment Designs}

\subsubsection*{EPO Dose-Response Sweep at Single-Cell Scale}

To characterise the EPO sensitivity of commitment probability at physiological
copy numbers, we ran $N=50$ independent stochastic replicates at each of eight
EPO concentrations spanning a 7\% range (0.430--0.460~mM). GCSF was held fixed
at 0.100~mM throughout. Each simulation ran for $T=7200$~s and was scored as
erythroid if GATA1\textsubscript{nuc}\,$>$\,1.5\,$\times$\,PU.1\textsubscript{nuc}
at $t=3600$~s. A Kruskal--Wallis test was used to assess whether GATA1\textsubscript{nuc}
at $t=3600$~s varied significantly across EPO conditions.

\subsubsection*{GCSF/EPO Ratio Sweep}

To map the full bifurcation structure of the two-cytokine control surface, we
ran $N=50$ stochastic replicates at each of ten GCSF concentrations
(0.5, 1.0, 1.2, 1.35, 1.5, 2.0, 2.2, 2.5, 2.8, and 3.0~mM) at EPO\,=\,1.0~mM
fixed. Each simulation ran for $T=7200$~s. Commitment scoring and Wilson score
95\% confidence intervals were applied as above.

\subsubsection*{Cytokine-Withdrawal Experiment}

A single deterministic trajectory was run at EPO\,=\,1.0~mM for
$t=0$--500~s, after which EPO\_external was set to zero via a timed environment
event and the trajectory was continued to $t=3600$~s. The erythroid state was
verified at $t=490$~s prior to withdrawal (GATA1$_\text{nuc}=10.91$~mM). The
post-withdrawal trajectory was assessed at 500~s intervals.

\subsubsection*{EPO Pulse-Duration Sweep}

EPO was applied at 1.0~mM and removed by timed environment events at eight
pulse durations: 60, 90, 120, 150, 300, 500, 1000, and 1800~s. Each
single-trajectory simulation continued to $t=3600$~s. Fate outcome was scored
as erythroid if GATA1/PU.1\,$>$\,1.5 at $t=3600$~s.

\subsubsection*{Sensitivity Analysis of the Priming-Crossover Time}

The phosphorylation rate constant governing the pGATA1 activation step was
rescaled by factors $\alpha \in \{0.50, 0.80, 1.00, 1.20, 1.50\}$, all other
parameters held constant. The priming-crossover time was identified as the first
$t$ at which the mean-trajectory pGATA1\textsubscript{nuc} exceeded
PU.1\textsubscript{nuc}.

\subsubsection*{EPO--pH Factorial Sweep}

To characterise the joint effect of EPO concentration and nuclear pH on
commitment probability, we ran a factorial design of eight EPO concentrations
(0.52, 0.54, 0.56, 0.57, 0.59, 0.61, 0.63, 0.65~mM) crossed with three nuclear
pH values (7.0, 7.5, 8.0) at fixed GCSF\,=\,1.1~mM. Each of the 24 conditions
used $N=100$ stochastic replicates and $T=21600$~s simulation time, giving a
total of 2400 replicates. Nuclear pH was set as a global environment parameter
affecting the inhibition constant via equation~(\ref{eq:pHhill}). Commitment was
scored as above; stochastic bistability was assessed by computing the bimodality
coefficient (BC\,=\,($\mu_3^2+1$)/($\mu_4+3(N-1)^2/((N-2)(N-3))$)) on the
log\,GATA1\textsubscript{nuc}/PU.1\textsubscript{nuc} distribution for each
condition; BC\,$>$\,0.555 was used as the bistability threshold. The EPO
commitment threshold EPO* at each pH was estimated by fitting a logistic function
P(ERY)\,=\,1/(1+exp($-k(\text{EPO}-\text{EPO}^*)$)) to the P(ERY)--EPO data for
each pH separately. Because all cells in this coarse sweep committed well before
$T=21600$~s (stochastic commitment is complete by $t\approx 300$--$500$~s near
the bifurcation), the extended simulation time does not affect outcome
classification but was retained for consistency.

\subsubsection*{Fine-Grid EPO Scans for EPO* Localisation}

To determine EPO* at each pH with high precision, three dedicated fine-grid EPO
scans were performed after the coarse factorial sweep. Each scan targeted a
narrow EPO window around the preliminary EPO* estimate from the coarse scan,
using $N=100$--$200$ stochastic replicates per condition and
$t_{\mathrm{end}}=2000$~s (sufficient to resolve all stochastic trajectories).
Binary-search narrowing across successive scan generations confirmed EPO* to
±0.003~mM or better at each pH. The resulting values (Table~\ref{tab:ph_threshold})
supersede the coarse-scan logistic-fit estimates for all downstream comparisons.

\subsubsection*{Attractor-Depth pH Experiment}

To isolate the effect of nuclear pH on committed-cell attractor depth, we ran
$N=50$ stochastic replicates at EPO\,=\,1.0~mM with GCSF\,=\,0.001~mM (effectively
zero competition) at each of three pH values (7.0, 7.5, 8.0). All 150 replicates
were verified to commit to the erythroid attractor (P\,(ERY)\,=\,100\%);
residual PU.1\textsubscript{nuc} at $t=3600$~s was reported as mean\,$\pm$\,SD
and maximum.

\subsection{Statistical Methods}

Commitment fractions are reported with Wilson score 95\% confidence intervals.
Comparisons across conditions were performed with the Kruskal--Wallis test for
non-parametric rank data; pairwise tests where noted used Mann--Whitney U with
Bonferroni correction. The bimodality coefficient was computed using the
third and fourth standardised moments of the log-ratio distribution. All
statistical analyses used Python 3.11 with pandas 2.x, NumPy 1.26.x, and
SciPy 1.12.x. No proprietary software was used.

\subsection{Energy Conservation Validation}

Model thermodynamic consistency was verified by confirming conservation of total
adenine nucleotide (ATP\,+\,ADP) and total guanine nucleotide (GTP\,+\,GDP)
across all simulations. These sums should remain constant because the model
contains no synthesis or degradation of nucleotide bases---only interconversion
between phosphorylation states. Deviations would indicate numerical integration
errors or stoichiometric inconsistencies. All simulations maintained conservation
to machine precision.
